\section{Implementation and Deployment Status}
\label{sec:implementation}

\subsection{Training Anchor}

The main implementation anchor is \texttt{weeklyresult/weekly\_drone\_2026w14/preprocess\_ext}. Its run configuration uses \texttt{base} teacher and student profiles, batch size $32$, learning rate $0.0001$, student training up to $50$ epochs, and teacher reuse from \texttt{saved\_models/weekly\_drone\_2026w14/baseline/teacher\_clean\_best.weights.h5}. The noisy setup mixes noise with probability $1.0$, samples training SNR from $-15$ dB to $-5$ dB, and evaluates at $-10$ dB.

The same configuration enables class-aware prosody augmentation for the student path. For the emergency class, the recorded pitch range is $2.0$ to $5.0$ and the gain range is $6.0$ dB to $10.0$ dB. The paper should describe this as a controlled safety-oriented augmentation setting, not as a general-purpose recipe.

\subsection{Host-Side TFLite Export}

The Phase 2 export artifacts show that the \texttt{w14/preprocess\_ext} full-integer TFLite model has input shape $[1,256,32,1]$ and int8 input and output tensors. The artifact summary records a full-integer flatbuffer size of $1893928$ bytes. A host-side TFLite runtime profile reports $165$ tensors and $104$ operators, with the largest visible runtime tensor at $524288$ bytes and shape $[1,256,32,64]$.

The runtime profile also reports anonymous temporary tensors, including a $1179648$ byte tensor with shape $[1,64,32,576]$. This is why the deployment discussion should focus on runtime arena pressure and operator working sets, not only on model-file size.

\subsection{ESP32/TFLM Constraint}

The board-side budget records a current full-integer TFLM arena requirement of $8225824$ bytes for the full local pipeline, while the target range is $4194304$ to $6291456$ bytes. Removing transport and JSON buffers is still useful engineering work, but the Phase 2 artifacts show that model activations and temporary tensors are the dominant deployment pressure.

The later \texttt{xiao\_bottleneck256} board trace improves the memory picture but exposes an operator constraint. It records a full-integer flatbuffer of $818496$ bytes, allocation in a $4194304$ byte arena, successful \texttt{AllocateTensors()}, and then a crash during \texttt{Invoke()} at \texttt{Conv2D Eval}. The failing temporal convolution has input dimensions $[1,1,32,256]$, filter dimensions $[256,1,31,1]$, and output dimensions $[1,1,32,256]$.

\subsection{Current Deployment Candidate}

The current candidate is \texttt{xiao\_bottleneck256\_tflm}. It keeps the Branchformer entry at $(32,256)$ and records $712067$ parameters, but replaces the grouped temporal \texttt{Conv1D} form with a depthwise form that exports as \texttt{DEPTHWISE\_CONV\_2D}. The handoff's local export check reports \texttt{CONV\_2D=6} and \texttt{DEPTHWISE\_CONV\_2D=1}.

The correct implementation claim is limited: the candidate addresses the known grouped-convolution incompatibility in local export, and a matching \texttt{w17/B\_small\_teacher\_student} run is close to the \texttt{w14} noisy-set anchor. The evidence does not yet support a completed ESP32/TFLM deployment claim.
